\documentclass[11pt]{article}
\usepackage{shared/hanzocover}
\usepackage{amsmath,amssymb,amsthm,booktabs,hyperref,graphicx}
\usepackage{xcolor}
\usepackage{listings}
\input{shared/lstlang}

\definecolor{codegreen}{rgb}{0,0.6,0}
\definecolor{codegray}{rgb}{0.5,0.5,0.5}
\definecolor{codepurple}{rgb}{0.58,0,0.82}
\definecolor{backcolour}{rgb}{0.95,0.95,0.92}

\lstdefinestyle{mystyle}{
    backgroundcolor=\color{backcolour},
    commentstyle=\color{codegreen},
    keywordstyle=\color{codepurple},
    numberstyle=\tiny\color{codegray},
    stringstyle=\color{codegreen},
    basicstyle=\ttfamily\footnotesize,
    breakatwhitespace=false,
    breaklines=true,
    captionpos=b,
    keepspaces=true,
    numbers=left,
    numbersep=5pt,
    showspaces=false,
    showstringspaces=false,
    showtabs=false,
    tabsize=2
}
\lstset{style=mystyle}

\newtheorem{theorem}{Theorem}
\newtheorem{lemma}{Lemma}
\newtheorem{proposition}{Proposition}
\newtheorem{definition}{Definition}
\newtheorem{corollary}{Corollary}

\title{Hanzo ML and Hanzo Engine: A Unified Multi-Backend, Multi-Modal Inference Stack}
\author{Hanzo Dev \\ \texttt{dev@hanzo.ai} \\ \textit{Hanzo AI}}
\date{2026}

\begin{document}

\hanzocoverpage

\begin{abstract}
We present Hanzo ML and Hanzo Engine, a single Rust inference stack that runs the same model
across five compute backends (CUDA, ROCm, Metal, Vulkan, CPU), serves six modalities (text,
vision, audio, speech, image generation, embeddings), and carries the entire GGUF quantization
zoo through one unified compute core. The thesis is structural: one templated quantization core,
one model loader, one set of pipelines, with each axis (weight format, activation dtype,
dense-versus-mixture-of-experts, scalar-versus-\texttt{dp4a}) orthogonal and composable. Adding a
quantization format costs one device decode function plus one traits row plus one table entry, and
no new kernel; we state this as a constant-cost extension proposition and report all $23$
native-resident formats decoding bit-exact ($\texttt{nbad}=0$) against the CPU reference. The
performance claims are measured head-to-head against \texttt{llama.cpp} (which is the engine inside
Ollama, so these are also the Ollama-class comparison) on the \emph{same} GPU, same model, same
quantization, under a device-map-controlled protocol. On the AMD Ryzen AI Max+ 395 ``Strix Halo''
APU (\texttt{gfx1151}) Hanzo \emph{beats} \texttt{llama.cpp} on prefill for every model tested
against the HIP backend ($1.02$--$1.14\times$ dense, $1.13\times$ on a $30$B mixture-of-experts)
and against the Vulkan backend on the $4$B/$8$B/$30$B models ($1.14$--$1.26\times$); on decode it
beats \texttt{llama.cpp}-HIP on the $30$B-A3B mixture-of-experts ($1.20\times$) and on the
\texttt{IQ3\_XXS} i-quant ($1.13\times$), and reaches HIP parity on the dense workhorses
($0.90$--$0.96\times$). On the NVIDIA GB10 ``Grace Blackwell'' it beats \texttt{llama.cpp}-CUDA on
small-model decode ($1.10\times$ on $0.6$B via a captured decode graph), reaches decode parity
on $8$B, and leads CUDA mixture-of-experts prefill ($1.87$--$2.1\times$ via an expert-grouped GEMM).
We are equally explicit about the active frontier, marked in-flight with the lever named
rather than dressed up: decode versus \texttt{llama.cpp}-Vulkan on the APU ($0.85$--$0.96\times$,
an attention-epilogue fusion); CUDA \emph{dense}-attention prefill, where a Blackwell-tuned
flash-attention kernel is now built, numerically stable on \texttt{sm\_121}, and shipped in the CUDA
images, lifting dense prefill $+42\%$ (\texttt{pp1024} $1425.5 \to 2030.9$ tok/s) and closing the gap
to \texttt{llama.cpp} from $0.55\times$ to $0.77$--$0.84\times$ yet leaving dense-attention prefill the
one axis Hanzo still trails; and the $35$B hybrid gated-delta-net prefill on CUDA, where a
warp-per-column recurrence now \emph{beats} \texttt{llama.cpp}'s GDN kernel ($0.93$ vs $1.04$
ms/layer, total prefill $354 \to 1166$ tok/s) with a \texttt{bf16} tensor-core GEMM the next lever for
the projections and experts. Every number here is measured on named silicon at a named engine version; where a path
is a projection or a microbenchmark it is labelled. The companion paper~\cite{hanzorocm} carries
the formal roofline analysis and the kernel-level campaign; this paper is the cross-backend,
cross-modality overview, and its honesty is the point: the data carries the story without
embellishment.
\end{abstract}

\section{Introduction}
\label{sec:intro}

Most inference stacks are a federation: a CUDA engine, a separate Metal port, a third thing for
AMD, a fourth for CPU, and per-format quantization kernels copied and specialized until the matrix
of (backend $\times$ format $\times$ dense/MoE) is unmaintainable. Hanzo takes the opposite
position, the one our engineering culture states as a law: one and only one way to do everything,
composable and orthogonal. The same Hanzo Engine build runs CUDA, ROCm, Metal, Vulkan, and CPU
behind one model loader; the same Hanzo ML quantization core serves the full 1-bit-to-8-bit GGUF
zoo across all of them; the same pipelines drive text, vision, audio, speech, image generation, and
embeddings. There is no per-vendor inference stack to maintain and no rewrite to move a workload
from a Mac to a Grace-Blackwell box to an APU mini-PC.

This paper makes two kinds of claims, and keeps them separate. The first is \emph{breadth}, which
is a matter of what the code does: the backend set, the modality set, the architecture set, and the
quant zoo (Sections~\ref{sec:ml} and~\ref{sec:engine}). The second is \emph{performance}, which is
a matter of measurement: same-GPU, same-quant, same-metric head-to-head numbers against
\texttt{llama.cpp} (Section~\ref{sec:results}). We treat \texttt{llama.cpp}~\cite{llamacpp,ggml} as
the baseline throughout because it is the de-facto standard for edge and commodity LLM inference,
and because Ollama~\cite{ollama} wraps \texttt{llama.cpp}/\texttt{ggml}: our same-GPU
\texttt{llama.cpp} numbers are therefore the Ollama-class comparison, stated once here so it need
not be repeated. We do not benchmark against vLLM~\cite{kwon2023}, whose home is batched datacenter
serving; this paper is single-stream throughput on commodity unified-memory hardware.

The campaign behind these numbers has an ethos worth stating up front, because it determines how
to read the results: \textbf{measured honesty}. The companion paper~\cite{hanzorocm} documents
negative results as first-class citizens, a fully-debugged HIP-graph effort that delivered zero
speedup, a flash-decode kernel built and reverted, a lossless speculative decoder that is a net
slowdown on this hardware. We continue that discipline here. Where Hanzo leads we show the data and
let it carry; where it trails we say so, quantify the gap, and name the one lever that closes it.
A paper of real measured wins plus honest frontier framing is more credible, and more useful to
the people deciding what to deploy, than a paper that claims to win everything.

\paragraph{Summary of measured standing.} Against \texttt{llama.cpp} on the same silicon:
\begin{itemize}
  \item \textbf{Prefill, AMD \texttt{gfx1151}:} Hanzo \emph{wins} on every model versus the HIP
  backend ($1.02$--$1.14\times$ dense; $1.13\times$ on the $30$B-A3B MoE) and versus the Vulkan
  backend on the $4$B/$8$B/$30$B models ($1.14$--$1.26\times$).
  \item \textbf{Decode, AMD \texttt{gfx1151}:} Hanzo \emph{wins} versus HIP on the $30$B-A3B MoE
  ($1.20\times$) and on \texttt{IQ3\_XXS} ($1.13\times$); reaches HIP parity on dense $4$B/$8$B
  ($0.90$--$0.96\times$); trails the Vulkan backend ($0.85$--$0.96\times$), the named frontier.
  \item \textbf{Decode, NVIDIA GB10:} Hanzo \emph{wins} versus CUDA on $0.6$B ($1.10\times$, via the
  captured decode graph) and reaches parity on $8$B.
  \item \textbf{Prefill, NVIDIA GB10:} Hanzo \emph{leads} CUDA MoE prefill ($1.87$--$2.1\times$ via the
  expert-grouped GEMM) and the GDN recurrence kernel ($0.93$ vs $1.04$ ms/layer; $35$B prefill
  $354 \to 1166$ tok/s). \emph{Dense}-attention prefill is the lone trailing axis: flash-attention is
  now shipped and lifts it $+42\%$ ($0.55\times \to 0.77$--$0.84\times$ \texttt{llama.cpp}), but it
  still trails at $\approx 0.8\times$.
  \item \textbf{Coverage:} the only portable multi-backend stack carrying all $22$ GGUF formats plus
  MXFP4 native-resident through a single compute core, every one bit-exact; full multimodality
  including a unified omni model (text, vision, audio, speech).
  \item \textbf{In-flight (lever named):} CUDA dense prefill residual ($0.77$--$0.84\times$ with
  flash-attention shipped; the remaining $\approx 20\%$ is the \texttt{f32}$\leftrightarrow$\texttt{bf16}
  cast chain plus MMQ GEMM efficiency); the $35$B hybrid total prefill (the GDN recurrence kernel now
  leads; \texttt{bf16} tensor-core GEMM next for the projections and experts); APU decode versus Vulkan
  ($0.96\times$; attention-epilogue fusion).
\end{itemize}

\section{Hanzo ML: The Unified Quantization Compute Core}
\label{sec:ml}

Hanzo ML is the tensor and kernel layer: a Rust framework with a device abstraction over five
backends and a quantization core that is the technical heart of the stack. Its design commitment,
a standing user mandate, is to support quantization from $1$-bit to full precision with \emph{no
per-quant kernels}: one quant-agnostic GPU core per stage, one decode function per format, reused
across decode, prefill, and mixture-of-experts.

\subsection{The add-a-quant contract and the orthogonal axes}
\label{sec:axes}

The core is templated on a weight-format type parameter \texttt{WTYPE} and factors the problem into
orthogonal axes that compose rather than multiply:
\begin{itemize}
  \item \textbf{Weight format} (\texttt{WTYPE}): the GGUF block layout and its dequantization, isolated
  in one device function \texttt{qdec<WTYPE>::partial} and one traits row \texttt{qdw\_traits<WTYPE>}
  (block bytes, elements, symmetry, sub-unit count).
  \item \textbf{Activation dtype} (\texttt{Act}): \texttt{f16}, \texttt{bf16}, or \texttt{f32}, selected at
  the entry point, so the residual stream can stay in its native dtype without a cast chain.
  \item \textbf{Dense versus MoE}: a single \texttt{bool MOE} template parameter; under
  \texttt{if constexpr(MOE)} the same core adds only an expert-grouped gather/scatter, and the dense
  path codegens byte-identical.
  \item \textbf{Scalar versus \texttt{dp4a}}: an int8 dot-product decode path
  (\texttt{qdp4a<WTYPE>}) sits beside the scalar path, trait-selected per type, with the scalar core
  as the always-correct fallback.
\end{itemize}
These axes are realized as three matched cores: \texttt{qmatvec\_core<WTYPE,Act>} for decode (with
its MoE twin and its \texttt{dp4a} twin), \texttt{qmmq\_core<WTYPE,MOE,NWAVE\_M>} for the int8
matrix-core prefill GEMM, and fused service kernels (\texttt{moe\_route} for softmax/top-$k$/normalize,
\texttt{moe\_combine} for the expert reduction, \texttt{rope\_norm\_positions} for fused
qk-RMSNorm-into-RoPE). The decomplection is deliberate: where three sibling efforts had each invented
a prefill-capability predicate, the merged core has exactly one, \texttt{RocmQuantType::qmmq\_capable},
read at every site. The format identifier is a \emph{value} (\texttt{RocmQuantType}) qualified by
namespace, not a place: dispatch reads the (type, activation) pair off the enum and selects the single
core's entry point. We state the resulting structural property precisely.

\begin{proposition}[Constant-cost format extension]
\label{prop:o1}
Let the unified core be parameterized by \texttt{WTYPE}. Adding support for a new quantization
format $F$ to the decode, prefill, and MoE paths requires exactly (i) one device decode function
\texttt{qdec<F>::partial} (and, optionally, one \texttt{qdp4a<F>} for the int8 path), (ii) one traits
row \texttt{qdw\_traits<F>}, and (iii) one dispatch-table entry \texttt{RocmQuantType::F}. No new
kernel, launcher, or dispatch branch is introduced; the matvec, GEMM, MoE gather/scatter, and
warp-reduction code are shared and unchanged. Hence the per-format marginal source is $O(1)$ and
independent of the number of already-supported formats.
\end{proposition}

\begin{proof}[Proof sketch]
By construction of the templated cores. The accumulation, reduction, and (for prefill) the
$128\times128$ \texttt{iu8}-WMMA tiling loop are written once and instantiated per \texttt{WTYPE}; the
only per-type code they call is \texttt{qdec<WTYPE>}/\texttt{qdp4a<WTYPE>} and the values in
\texttt{qdw\_traits<WTYPE>}. Dispatch is by symbolic kernel name keyed on \texttt{RocmQuantType}, so a
new variant adds one match arm, not a branch in the hot loop. The empirical witness is that the
six originally-wired formats and the sixteen added in one wave produce no new kernel and pass the
same numeric oracle; the refactor of the prefill GEMM onto the template produced byte-identical
disassembly. The formal bit-exactness of each instantiation is Lemma~1 of the companion
paper~\cite{hanzorocm}.
\end{proof}

\subsection{The complete quant zoo}
\label{sec:zoo}

The contract of Proposition~\ref{prop:o1} is realized to completion: \textbf{all $22$ GGUF
quantization formats plus MXFP4 ($23$ types total) are wired native-resident}, every one bit-exact
against the CPU \texttt{to\_float} reference ($\texttt{nbad}=0$ across $233$ oracle assertions, zero
regression on the pre-existing set). The types span every structural class (Listing~\ref{lst:zoo}).

\begin{lstlisting}[language=Rust,caption={The $23$ native-resident formats, one \texttt{qdec} body and one traits row each.},label={lst:zoo}]
// k-quants:   Q2_K Q3_K Q4_K Q5_K Q6_K          (256-elem super-blocks)
// legacy:     Q4_0 Q4_1 Q5_0 Q5_1 Q8_0 Q8_1     (32-elem blocks)
// i-quants:   IQ1_S IQ1_M IQ2_XXS IQ2_XS IQ2_S
//             IQ3_XXS IQ3_S IQ4_NL IQ4_XS        (codebook / grid LUT)
// ternary:    TQ1_0 TQ2_0                        (base-3 packed)
// fp4:        MXFP4                              (E8M0 scale + FP4 codebook, gpt-oss)
\end{lstlisting}

One \texttt{qmmq\_capable} predicate partitions the set: $11$ formats carry an int8-WMMA prefill
kernel; the rest dequantize to \texttt{f16} for prefill while decoding native. The whole zoo collapses
onto \emph{two} scalar accumulation shapes folded inside each decode (symmetric
$\text{scale}\cdot q$, asymmetric $d\cdot sc\cdot q - d_{\min}\cdot m$), so the cores themselves are
shape-agnostic. The honest distinction from \texttt{llama.cpp} is not coverage; \texttt{llama.cpp}
defines the GGUF formats and supports them all. It is the \emph{architecture}: Hanzo serves the
entire zoo through one \texttt{WTYPE}-templated core family, the same cores driving decode, prefill,
and MoE across five backends, with the one-line add-a-quant contract of
Proposition~\ref{prop:o1}; \texttt{llama.cpp} uses per-format front-ends and is primarily a single
vendor's kernel set per backend. Among portable, pure-Rust, multi-backend engines, Hanzo ML is the
only one carrying the full zoo native-resident through a single compute core.

\subsection{Five backends, one loader}
\label{sec:backends}

The \texttt{Device} abstraction exposes \texttt{Cpu}, \texttt{Cuda}, \texttt{Metal}, \texttt{Rocm},
and \texttt{Vulkan}, each with a full backend implementation behind one storage trait. A model is
loaded once and runs on whichever backend is present; the cross-backend decode parity reported in
Section~\ref{sec:results} is itself evidence that the unified core is not over-fit to one vendor.
What is \emph{ported} and what is \emph{novel} is stated plainly, as in the companion work: the
\texttt{mul\_mat\_q} integer-matrix-core prefill design and the GGUF block formats follow
\texttt{llama.cpp}/\texttt{ggml}; the unification (one \texttt{WTYPE}-templated core spanning matvec,
GEMM, and MoE with a one-line add-a-quant contract) is Hanzo's.

\subsection{Decode fusions and the bandwidth wall}
\label{sec:fusions}

The optimization program is organized by a single roofline dichotomy, which we state as the
operating principle of the whole stack and whose formal ceiling is proved in the
companion~\cite{hanzorocm,williams2009}.

\begin{proposition}[Two-regime roofline]
\label{prop:tworegime}
Single-stream decode (batch $1$, one new token) streams every weight once per token against a single
activation row; it is therefore \emph{memory-bandwidth bound}, with throughput at most $B/W$ tokens
per second for sustained bandwidth $B$ and resident footprint $W$. Prefill (a multi-token prompt)
multiplies a weight matrix by an activation matrix; it is \emph{compute bound} on the matrix or
tensor cores. Consequently decode optimization must reduce bytes streamed or kernel/launch overhead,
and prefill optimization must improve matrix-core scheduling; a technique aimed at the wrong regime
cannot move the bottleneck.
\end{proposition}

This dichotomy explains every decode result. On \texttt{gfx1151} the per-expert decode matvec already
sustains $88$--$99\%$ of the $212$ GB/s memory roofline (the Q6\_K matvec does \emph{more} arithmetic
yet reaches \emph{higher} bandwidth than Q4\_K, which is impossible under a compute bound), so the
matvec is finished and the remaining decode lever is \emph{kernel count}, not kernel speed. Acting on
that diagnosis, a sequence of fusions, each bit-exact and each a true elimination of work rather than
a mere launch merge, moved end-to-end decode: a fused \texttt{moe\_route} kernel collapsing the
six-launch-per-layer routing chain to one ($+13\%$ on the $30$B-A3B MoE); a fused gate/up shared
activation-quantize plus a fused residual-add RMSNorm ($+2.8\%$); a fused qk-RMSNorm-into-RoPE kernel
($+2.8\%$). The i-quants, which decoded through a scalar codebook lookup and were pathologically
compute-bound, were given a \texttt{dp4a} grid-dot path that lifted \texttt{IQ3\_XXS} decode
$2.27\times$ (from $30.8$ to $\approx 68$ tok/s), enough to \emph{pass} \texttt{llama.cpp}-HIP on that
format. These are the wins behind the decode numbers in Section~\ref{sec:results}.

\subsection{Honest negatives, as first-class results}
\label{sec:negatives}

The same measurement discipline ruled out a set of plausible levers, and recording them is part of
the contribution because each one would otherwise be re-attempted.
\begin{itemize}
  \item \textbf{RDNA3.5 \texttt{iu8} WMMA equals \texttt{f16} WMMA ($1.00\times$).} A register-only
  microbench measured $57.1$ TOP/s int8 versus $56.8$ TFLOP/s \texttt{f16}; unlike CDNA's $2\times$
  int8 advantage, RDNA3.5 has none. So on this APU quantization is a decode-side bandwidth/footprint
  play and \emph{never} a prefill-side compute play. A hardware limit, not a tuning gap.
  \item \textbf{HIP-graph decode is correct but inert when already de-launch-bounded.} Capturing and
  replaying the per-token decode graph delivered exactly zero speedup once the fusion work had removed
  the launches it was built to collapse, reproduced at two scales ($8$B and a $30$B MoE) and two
  platforms. Graphs pay only where launch-bound (small models, and the MoE routing chain before it was
  fused), which is exactly where we now enable them.
  \item \textbf{Matvec rewrites are flat at the wall.} \texttt{f32}-native activation streaming, a
  fused router GEMV, and dense lane-coalescing are each bit-exact yet move device-map-controlled decode
  by $\leq 1.3\%$, because they optimize a kernel already at $88$--$99\%$ of bandwidth. Folding the
  activation quantize \emph{into} the matvec was bit-exact but $-22\%$: it injects redundant per-block
  work into a bandwidth-bound kernel. Reverted.
  \item \textbf{Lossless speculative decode is a net slowdown here.} A classic draft-plus-target
  decoder, added as a second proposer on the existing trait, is provably lossless at temperature $0$
  (greedy output token-identical with and without a draft) yet measures $0.4$--$0.63\times$ on this APU
  with stock small drafts: acceptance is only $33$--$40\%$, the verify is charged at the prefill rate,
  and a $3$B-active MoE has no per-token headroom to amortize. Shipped, proven, and recommended against
  here; the multi-token-prediction path on the same trait is the viable variant.
\end{itemize}
The pattern is consistent: we profile the production kernel, attribute every millisecond, optimize
only what the profiler proves is the bottleneck, and report the disproven ideas as valid negative
results.

\section{Hanzo Engine: Architectures, Modalities, Pipelines}
\label{sec:engine}

Hanzo Engine is the inference server and model layer on top of Hanzo ML: the model implementations,
the pipelines (text, multimodal, speech, diffusion, embedding), the OpenAI- and Anthropic-compatible
HTTP surface, the MCP and Python and Rust SDKs, PagedAttention, and the decode-graph and speculative
machinery. Its breadth is the second half of the thesis.

\subsection{Architecture breadth}
\label{sec:arch}

The engine ships $32$ text architectures and $27$ vision modules behind one loader. The set spans
three structural families that matter for inference, summarized in Table~\ref{tab:arch}.
\begin{itemize}
  \item \textbf{Dense transformers:} Llama, Mistral, Qwen2/Qwen3, Gemma/Gemma2/Gemma3, Phi2/Phi3,
  Granite, Starcoder2, SmolLM3.
  \item \textbf{Mixture-of-experts:} Qwen3-MoE (the $30$B-A3B workhorse), Qwen3.5-MoE, Mixtral,
  Phi3.5-MoE, GLM4-MoE, GPT-OSS (MXFP4 experts), and DeepSeek-V2/V3 (multi-head latent attention
  plus MoE).
  \item \textbf{Hybrid linear-attention plus MoE:} Qwen3-Next, which interleaves gated-delta-net
  (GDN) linear-attention layers~\cite{gateddeltanet} with a full-attention layer every fourth layer
  and an MoE feed-forward. This is the $35$B-A3B hybrid benchmarked in Section~\ref{sec:cuda}; it
  exercises a recurrence (the chunked gated delta rule) and a causal conv1d in addition to the usual
  attention and MoE kernels.
\end{itemize}

\begin{table}[t]
\centering
\caption{Architecture and modality breadth (representative, not exhaustive).}
\label{tab:arch}
\begin{tabular}{ll}
\toprule
\textbf{Family / modality} & \textbf{Representative models} \\
\midrule
Dense LLM        & Llama, Mistral, Qwen3, Gemma3, Phi3, Granite, Starcoder2, SmolLM3 \\
MoE LLM          & Qwen3-MoE, Qwen3.5-MoE, Mixtral, Phi3.5-MoE, GLM4-MoE, GPT-OSS, DeepSeek-V3 \\
Hybrid GDN+MoE   & Qwen3-Next (gated-delta-net linear attn + full attn + MoE) \\
Vision-language  & Qwen2.5-VL, Qwen3-VL(-MoE), Llama4, Gemma3/3n, Idefics2/3, LLaVA, Mistral3, Phi4 \\
Audio in         & Voxtral, Conformer, Qwen3-ASR, MiniCPM-O \\
Speech out (TTS) & Dia, Qwen3-TTS \\
Omni             & Qwen3-Omni (text + vision + audio + speech) \\
Image generation & FLUX, Qwen-Image (plus MuseTalk lip-sync) \\
Embeddings       & Qwen3-Embedding, EmbeddingGemma \\
\bottomrule
\end{tabular}
\end{table}

\subsection{Modality breadth}
\label{sec:modality}

The pipeline layer exposes six model categories: text, vision, audio, speech, diffusion (image
generation), and embeddings. Vision-language support is broad (Qwen2.5-VL, the Qwen3-VL family,
Llama4, the Gemma3 line, Idefics2/3, LLaVA, Mistral3, Phi4, with CLIP/SigLIP encoders); audio input
runs through Voxtral, a Conformer, and Qwen3-ASR; speech synthesis runs through Dia and Qwen3-TTS;
image generation runs through FLUX and Qwen-Image. The capstone is \textbf{Qwen3-Omni}, a single
model that ingests text, images, and audio and emits text and speech, served by the same engine and
the same quantized cores as the text models. This is the multimodal breadth the stack claims: not a
bag of separate tools, but one engine and one quantization core spanning the modalities.

\subsection{Pipelines: paged attention, decode graphs, speculative decode}
\label{sec:pipelines}

Three pieces of pipeline machinery are worth naming because they recur in the results.
\textbf{PagedAttention}~\cite{kwon2023} is implemented on every backend (the ROCm port was a
keystone built during the campaign and is the prerequisite for the speculative and prefix-cache
paths). \textbf{Decode graphs} capture the steady-state single-token decode and replay it as one GPU
graph instead of relaunching hundreds of kernels per token; they are wired for safetensors and GGUF
on both CUDA and ROCm, gated by an eligibility invariant (only architectures whose decode RoPE reads
a stable device position buffer are graph-safe, which excludes host-offset and per-forward mRoPE
variants that would otherwise freeze the position at capture and emit garbage on replay). The CUDA
GGUF decode graph is the lever behind the small-model CUDA decode win in Section~\ref{sec:cuda}.
\textbf{Speculative decoding} is factored behind one \texttt{SpeculativeProposer} trait with two
implementors, draft-model and multi-token-prediction, composed onto the same verifier; the draft path
is proved lossless at temperature $0$ (Section~\ref{sec:negatives}).

\section{Method}
\label{sec:method}

Every performance number in this paper is measured on named silicon (Table~\ref{tab:hw}), at a named
engine version, against \texttt{llama.cpp} on the \emph{same} GPU, same model, same quantization, quiet GPU. The
protocol controls the confounds that make unified-memory benchmarking treacherous.

\begin{itemize}
  \item \textbf{Same-GPU, same-quant, same-metric.} Prefill is $512$- or $1024$-token prompt
  throughput (\texttt{pp512}/\texttt{pp1024}); decode is single-stream token generation at shallow KV
  depth. For \texttt{llama.cpp} these are \texttt{llama-bench}'s \texttt{pp}/\texttt{tg}; for Hanzo they
  are the \texttt{hanzo bench} prefill and decode at matched lengths. We compare against the
  \texttt{llama.cpp} HIP \emph{and} Vulkan backends on AMD, because \texttt{llama.cpp} itself decodes
  faster on Vulkan than HIP on this APU, and comparing only against its slower backend would flatter us.
  \item \textbf{Device-map control.} On a unified-memory box the auto device-mapper can misread shared
  RAM as per-layer VRAM and silently offload roughly a third of the layers to the CPU, which collapses
  decode (prefill, being compute-bound, survives) and manufactures illusory speedups. We force all
  layers resident on the accelerator for both engines (\texttt{-ngl 99} for \texttt{llama.cpp};
  \texttt{-n 0:NLAYERS} for Hanzo). An uncontrolled A/B against a mis-mapped baseline once produced an
  apparent ``$+28\%$'' that was entirely the device fix; held constant, that same change was flat.
  \item \textbf{Thermal control.} The APU's warm runs swing $\approx 20\%$, so sequential comparison is
  untrustworthy; A/B comparisons are thermally interleaved and ratios are taken per round.
  \item \textbf{Correctness gating.} Every kernel change is gated bit-exact ($\texttt{nbad}=0$) against
  a CPU numeric oracle before any timing is trusted, and output coherence is verified.
\end{itemize}

\begin{table}[t]
\centering
\caption{Hardware. All three are commodity unified-memory boxes (no discrete VRAM tier).}
\label{tab:hw}
\begin{tabular}{lll}
\toprule
\textbf{Box} & \textbf{Accelerator} & \textbf{Memory / stack} \\
\midrule
AMD Strix Halo & Radeon 8060S iGPU, \texttt{gfx1151} (RDNA3.5), 40 CU @ 2.9 GHz & 62 GB unified LPDDR5X, $\approx$231 GB/s; ROCm 7.13 (Linux) \\
NVIDIA GB10    & Grace-Blackwell, \texttt{sm\_121}, integrated GPU & 122 GB unified LPDDR5X; CUDA 13, Linux aarch64 \\
Apple M4 Max   & M4 Max GPU, Metal & 137 GB unified; \texttt{llama.cpp} \texttt{BLAS,MTL} \\
\bottomrule
\end{tabular}
\end{table}

\section{Results}
\label{sec:results}

\subsection{AMD \texttt{gfx1151}: prefill won, decode at HIP parity-or-better}
\label{sec:amd}

Table~\ref{tab:amd} is the headline same-GPU comparison on the Strix Halo APU: Hanzo
(ROCm~\cite{rocm2024}) against \texttt{llama.cpp} on both its HIP and its Vulkan backend, across three
dense Qwen3~\cite{qwen3} models and the $30$B-A3B mixture-of-experts. Two facts dominate.

First, \textbf{Hanzo wins prefill outright}. Against the HIP backend it is faster on every model
($1.02$--$1.14\times$ dense, $1.13\times$ on the $30$B MoE); against the heavily community-tuned
Vulkan backend it wins on the $4$B, $8$B, and $30$B models ($1.14$--$1.26\times$) and is within
$4\%$ on the tiny $0.6$B. The MoE prefill win is the headline kernel result of the campaign: a
fused expert-grouped WMMA GEMM plus a fused expert-combine took MoE prefill from $0.34\times$ to a
$1.13\times$ \emph{lead} over \texttt{llama.cpp}-HIP, entirely through kernels the engine owns.

Second, \textbf{decode is at HIP parity-or-better}. Hanzo beats \texttt{llama.cpp}-HIP on the
$30$B-A3B MoE ($1.20\times$, the production path with HIP-graphs default-on and the full decode-fusion
campaign) and reaches HIP parity on the $4$B/$8$B workhorses ($0.90$--$0.96\times$). It trails HIP
only on the tiny $0.6$B ($0.69\times$), where decode is dispatch-overhead-bound and the fixed
per-token launch cost dominates a small weight stream. Against the Vulkan backend, decode trails by
$0.85$--$0.96\times$: this is the active frontier (Section~\ref{sec:frontier}), and the gap is
structural, Vulkan's sustained occupancy and fused attention epilogue, not a slower matvec (ours is
at the bandwidth wall).

\begin{table}[t]
\centering
\caption{AMD \texttt{gfx1151}, Hanzo ROCm vs \texttt{llama.cpp} HIP and Vulkan on the \emph{same}
GPU (\texttt{pp1024}/\texttt{tg128}; the $30$B-A3B is Q4\_K\_M, decode on the production HIP-graph
path). Throughput in tok/s; ratio is Hanzo / \texttt{llama.cpp}. Bold marks a Hanzo lead.}
\label{tab:amd}
\begin{tabular}{l rrr rrr}
\toprule
& \multicolumn{3}{c}{\textbf{Prefill (tok/s)}} & \multicolumn{3}{c}{\textbf{Decode (tok/s)}} \\
\cmidrule(lr){2-4}\cmidrule(lr){5-7}
\textbf{Model} & \textbf{Hanzo} & \textbf{vs HIP} & \textbf{vs Vk} & \textbf{Hanzo} & \textbf{vs HIP} & \textbf{vs Vk} \\
\midrule
Qwen3-0.6B-Q8\_0 (dense)   & 6966 & \textbf{1.14$\times$} & 0.96$\times$           & 128  & 0.69$\times$          & 0.60$\times$ \\
Qwen3-4B-Q8\_0 (dense)     & 1750 & \textbf{1.02$\times$} & \textbf{1.19$\times$}  & 58.2 & 0.96$\times$          & 0.85$\times$ \\
Qwen3-8B-Q8\_0 (dense)     & 1066 & \textbf{1.05$\times$} & \textbf{1.14$\times$}  & 34.6 & 0.90$\times$          & 0.87$\times$ \\
Qwen3-30B-A3B (MoE)        & 1209 & \textbf{1.13$\times$} & \textbf{1.26$\times$}  & 79.5 & \textbf{1.20$\times$} & 0.96$\times$ \\
\bottomrule
\end{tabular}
\end{table}

\subsection{The per-quant decode ladder}
\label{sec:ladder}

Because the quant zoo is the architectural centerpiece, Table~\ref{tab:ladder} shows decode across
three formats on the $30$B-A3B MoE in one controlled cross-quant pass (the \texttt{@d4} bench harness
at an intermediate engine version, used here for the relative comparison; the production decode-fusion
path of Table~\ref{tab:amd} raises the absolute Q4\_K figure further, as noted). Two readings matter.
The \textbf{k-quant bandwidth ladder works}: lower bits decode faster (Q4\_K $\to$ Q3\_K $\to$ Q2\_K),
consistent with Proposition~\ref{prop:tworegime}, and Hanzo is at HIP parity across the k-quants. And
\textbf{the i-quant \texttt{dp4a} path is a genuine win}: after vectorizing the codebook decode,
\texttt{IQ3\_XXS} decode rose $2.27\times$ and now \emph{beats} \texttt{llama.cpp}-HIP ($1.13\times$).
The residual gap to the Vulkan backend is uniform across formats and is the same structural
attention-epilogue gap as in Table~\ref{tab:amd}, not a per-format kernel deficit.

\begin{table}[t]
\centering
\caption{AMD \texttt{gfx1151}, Qwen3-30B-A3B decode by quantization format (cross-quant
\texttt{@d4} pass). Throughput in tok/s on the same GPU; ratio is Hanzo / \texttt{llama.cpp}.}
\label{tab:ladder}
\begin{tabular}{l r rr rr}
\toprule
\textbf{Format} & \textbf{bpw} & \textbf{Hanzo} & \textbf{llama-HIP} & \textbf{llama-Vk} & \textbf{vs HIP} \\
\midrule
Q4\_K   & 4.5  & 64.7 & 66 & 83 & 0.98$\times$ \\
Q2\_K   & 2.6  & 75.0 & 79 & 92 & 0.95$\times$ \\
IQ3\_XXS & 3.06 & 65.5 & 58 & 91 & \textbf{1.13$\times$} \\
\bottomrule
\end{tabular}
\end{table}

\subsection{NVIDIA GB10: decode and MoE prefill won, dense prefill the frontier}
\label{sec:cuda}

On the Grace-Blackwell GB10, Table~\ref{tab:cuda} compares Hanzo (CUDA) against
\texttt{llama.cpp}-CUDA on the same GPU and GGUF. The decisive recent result is the \textbf{CUDA GGUF
decode graph}: wiring the captured decode graph into the GGUF pipeline (it had been running eager,
$\approx 780$ kernel launches per token) lifted $0.6$B decode from $175.6$ to $247.5$ tok/s ($+41\%$),
crossing \texttt{llama.cpp}-CUDA ($225.6$) to $1.10\times$, byte-identical to the eager output. On the
$8$B, decode is at parity ($\approx 0.98\times$): there the model is bandwidth-bound and the graph adds
little, exactly as Proposition~\ref{prop:tworegime} predicts.

Prefill on CUDA splits cleanly by regime, and the split is the honest story. On the \emph{sparse} and
\emph{recurrent} paths Hanzo leads: the expert-grouped MoE prefill GEMM ported to CUDA (the same kernel
family that took AMD MoE prefill to a $1.13\times$ lead, Section~\ref{sec:amd}) delivers
$1.87$--$2.1\times$ MoE prefill over \texttt{llama.cpp}-CUDA, and on the $35$B-A3B hybrid GDN+MoE a
warp-per-column gated-delta-rule recurrence now \emph{beats} \texttt{llama.cpp}'s GDN kernel ($0.93$ vs
$1.04$ ms/layer). That hybrid shows the full campaign trajectory: de-redundifying the chunked recurrence
(eliminating $\approx 370\times$ redundant \texttt{exp} calls, $18.3 \to 4.0$ ms/layer), wiring conv1d to
a CUDA kernel, fusing the gated RMSNorm, and finally the warp-per-column recurrence took its prefill from
$354$ to $1166$ tok/s (a $3.3\times$ gain over the floor), bit-exact to an \texttt{f64}-truth reference at
$\approx 10^{-7}$ and coherent (the silent CUDA MoE routing-corruption bug that once panicked it is fixed
and gated by a regression oracle). The residual on the $35$B \emph{total} prefill is the projection and
expert GEMM; the named next lever is a \texttt{bf16} tensor-core GEMM there.

The one axis where Hanzo still trails is \emph{dense}-attention prefill, and Section~\ref{sec:flashattn}
reports it without inflation. Profiling attributed the GGUF dense-prefill gap ($0.33$--$0.55\times$) to
un-fused attention (softmax plus a materialized score tensor) and \texttt{f32}$\leftrightarrow$\texttt{bf16}
casts, with the actual MMQ GEMM only a small slice. The named lever, a Blackwell-tuned
flash-attention~\cite{dao2022} kernel, is now \emph{built and shipped} (Table~\ref{tab:flashattn}): it
lifts $8$B dense prefill $+42\%$ and closes the gap to $0.77$--$0.84\times$, but does not erase it.

\begin{table}[t]
\centering
\caption{NVIDIA GB10, Hanzo CUDA vs \texttt{llama.cpp}-CUDA on the \emph{same} GPU and GGUF.
Prefill/decode are Hanzo $/$ \texttt{llama.cpp} ratios; the $8$B prefill cell is no-FA$\to$flash-attn
(full prompt-length sweep in Table~\ref{tab:flashattn}), and the $35$B prefill cell is throughput
trajectory in tok/s (its GDN recurrence kernel leads \texttt{llama.cpp} at $0.93$ vs $1.04$ ms/layer).
Decode in tok/s where shown. Bold marks a Hanzo lead.}
\label{tab:cuda}
\begin{tabular}{l ll l}
\toprule
\textbf{Model} & \textbf{Prefill} & \textbf{Decode} & \textbf{Status / lever} \\
\midrule
Qwen3-0.6B (dense) & 0.33$\times$ & \textbf{1.10$\times$} (247.5 vs 225.6) & decode-graph shipped; dense prefill via flash-attn (Table~\ref{tab:flashattn}) \\
Qwen3-8B (dense)   & 0.55$\to$0.80$\times$ & $\approx$0.98$\times$ & flash-attn shipped ($+42\%$); residual: casts $+$ MMQ GEMM \\
Qwen3-Next-35B (GDN+MoE) & $354\to1166$ & runs coherent & GDN kernel leads ($0.93$ vs $1.04$ ms/layer); \texttt{bf16} GEMM next \\
\bottomrule
\end{tabular}
\end{table}

\subsection{Flash-attention on Blackwell: the dense-prefill lever, shipped}
\label{sec:flashattn}

Dense-attention prefill is the single axis on which Hanzo trails \texttt{llama.cpp}, and the named lever
to close it was a Blackwell-tuned flash-attention~\cite{dao2022} kernel. That kernel now exists, is
correct, and ships. We report the result without inflation (Table~\ref{tab:flashattn}): on the GB10
(\texttt{sm\_121}, CUDA $13.0$), same Qwen3-8B Q4\_K\_M GGUF as \texttt{llama.cpp}-CUDA, flash-attention
lifts dense prefill \textbf{$+42\%$} at \texttt{pp1024} ($1425.5 \to 2030.9$ tok/s) and holds the gain
across context lengths. Crossing the kernel off as built is itself a result: the CUTLASS-for-Blackwell
attention path compiles and runs \emph{numerically stable} on \texttt{sm\_121} (coherent generation, no
overflow), clearing the real risk that the new architecture would demand a different kernel.

The honest reading is the point. Flash-attention closes the dense-prefill gap from $0.55\times$ (no-FA)
to $\mathbf{0.77}$--$\mathbf{0.84\times}$ \texttt{llama.cpp}, a large improvement, \emph{but
dense-attention prefill remains the one axis where Hanzo still trails}, at $\approx 0.8\times$. We do not
claim a dense-prefill win. The residual $\approx 16$--$23\%$ is the part of the prefill path that
flash-attention does not touch: the \texttt{f32}$\leftrightarrow$\texttt{bf16} cast chain around the
quantized GEMM (measured $\approx 20\%$) and MMQ GEMM efficiency. Those are the next two levers, and
until they land this axis is stated as trailing, not won.

The win is shipped, not a branch. Flash-attention is in the CUDA container images:
\texttt{Dockerfile.cuda} builds \texttt{WITH\_FEATURES="cuda,cudnn,flash-attn"} on a CUDA $13.0.3$
base, and the build matrix now includes Blackwell \texttt{sm\_120}/\texttt{sm\_121}, so GB10 and
RTX-5090-class hardware get first-class CUDA images and the $+42\%$ reaches users rather than sitting in
a local build.

\begin{table}[t]
\centering
\caption{NVIDIA GB10 (\texttt{sm\_121}, CUDA $13.0$), Qwen3-8B Q4\_K\_M GGUF dense prefill: Hanzo CUDA
(without and with flash-attention) vs \texttt{llama.cpp}-CUDA on the \emph{same} GPU and GGUF.
Throughput in tok/s (higher is better); ``FA vs llama'' is Hanzo$+$flash-attn $/$ \texttt{llama.cpp}.
Flash-attention lifts dense prefill $+42\%$ and closes the gap to $0.77$--$0.84\times$, but dense
prefill remains the one axis Hanzo trails.}
\label{tab:flashattn}
\begin{tabular}{l rrr r}
\toprule
\textbf{Prompt} & \textbf{Hanzo no-FA} & \textbf{Hanzo $+$FA} & \textbf{\texttt{llama.cpp}} & \textbf{FA vs llama} \\
\midrule
\texttt{pp1024} & $1425.5$ & $2030.9$ & $2542.4$ & $0.80\times$ \\
\texttt{pp2048} & ---      & $2126.4$ & $2537.1$ & $0.84\times$ \\
\texttt{pp4096} & ---      & $1902.2$ & $2465.5$ & $0.77\times$ \\
\bottomrule
\end{tabular}
\end{table}

\subsection{Cross-backend: decode parity everywhere, and the box flips with the workload}
\label{sec:crossbackend}

A separate measurement pass~\cite{hanzorocm} fixed one model (Qwen3-8B-Q8\_0) and compared Hanzo to
\texttt{llama.cpp} across AMD, NVIDIA, and Apple (Table~\ref{tab:cross}). Two cross-cutting facts
fall out. First, \textbf{Hanzo reaches decode parity with \texttt{llama.cpp} on all three backends}
($0.98\times$/$1.04\times$/$0.99\times$ on AMD/NVIDIA/Apple), evidence that the unified core streams
at the hardware ceiling regardless of vendor; prefill trails by a backend-dependent
$0.76$--$0.91\times$ on that older pass (the AMD prefill has since become a lead, Table~\ref{tab:amd}).
Second, and orthogonal to the engine, \textbf{the fastest commodity box depends on the workload}:
NVIDIA GB10 wins prefill (compute-bound; Blackwell tensor cores), Apple M4 Max wins decode by roughly
$2.5\times$ (bandwidth-bound; widest unified memory), and the GB10, the strongest prefill box, is the
slowest at decode. A RAG or document-ingest workload prefers the GB10; an interactive chat or
long-generation workload prefers the M4 Max. The practical payoff of the unified stack is that the
\emph{same} build routes to whichever box the workload wants, with no per-vendor rewrite.

\begin{table}[t]
\centering
\caption{Cross-backend, Qwen3-8B-Q8\_0, single GPU (companion measurement pass~\cite{hanzorocm}).
Decode parity on all three; the hardware optimum flips between prefill and decode.}
\label{tab:cross}
\begin{tabular}{l rr c rr c}
\toprule
& \multicolumn{3}{c}{\textbf{Prefill (tok/s)}} & \multicolumn{3}{c}{\textbf{Decode (tok/s)}} \\
\cmidrule(lr){2-4}\cmidrule(lr){5-7}
\textbf{Backend} & \textbf{llama} & \textbf{Hanzo} & \textbf{ratio} & \textbf{llama} & \textbf{Hanzo} & \textbf{ratio} \\
\midrule
AMD Strix Halo (ROCm)  & 1176 & 1012 & 0.86$\times$ & 25.0 & 24.5 & 0.98$\times$ \\
NVIDIA GB10 (CUDA)     & 2155 & 1630 & 0.76$\times$ & 19.6 & 20.3 & 1.04$\times$ \\
Apple M4 Max (Metal)   &  873 &  792 & 0.91$\times$ & 50.4 & 49.9 & 0.99$\times$ \\
\bottomrule
\end{tabular}
\end{table}

\section{Where Hanzo Leads, and the Active Frontier}
\label{sec:frontier}

\paragraph{Where Hanzo leads (measured).}
\begin{enumerate}
  \item \textbf{Prefill on AMD \texttt{gfx1151}:} faster than \texttt{llama.cpp}-HIP on every model
  ($1.02$--$1.14\times$ dense, $1.13\times$ MoE) and faster than \texttt{llama.cpp}-Vulkan on
  $4$B/$8$B/$30$B ($1.14$--$1.26\times$). The MoE-prefill lead is the campaign's signature kernel win
  ($0.34\times \to 1.13\times$).
  \item \textbf{Sparse and recurrent prefill on NVIDIA GB10 (CUDA):} the expert-grouped MoE prefill
  GEMM leads \texttt{llama.cpp}-CUDA ($1.87$--$2.1\times$), and the warp-per-column GDN recurrence kernel
  beats \texttt{llama.cpp}'s ($0.93$ vs $1.04$ ms/layer; $35$B prefill $354 \to 1166$ tok/s).
  \emph{Dense}-attention prefill is the lone CUDA prefill axis still trailing (frontier item~2).
  \item \textbf{Decode wins where the regime allows:} $30$B-A3B MoE $1.20\times$ HIP; \texttt{IQ3\_XXS}
  $1.13\times$ HIP; $0.6$B $1.10\times$ CUDA via the decode graph; HIP and cross-backend decode parity
  on the dense workhorses.
  \item \textbf{Completeness:} all $22$ GGUF formats plus MXFP4 native-resident through one compute
  core, every one bit-exact; one engine across five backends; the full modality set including a unified
  omni model. This is the breadth no single-vendor or single-format stack matches.
\end{enumerate}

\paragraph{The active frontier (in-flight, lever named).} We mark these in-progress, not as losses,
because each has a profiler-identified lever and measured forward motion.
\begin{enumerate}
  \item \textbf{APU decode versus \texttt{llama.cpp}-Vulkan ($0.85$--$0.96\times$).} The matvec is at
  the bandwidth wall; the residual is Vulkan's fused attention epilogue and sustained occupancy. The
  lever is fusing the whole attention epilogue (norm, RoPE, cache-append, cast-free), a multi-kernel
  rewrite, not a single op.
  \item \textbf{CUDA dense-attention prefill ($0.77$--$0.84\times$, flash-attention shipped).} The
  Blackwell flash-attention kernel is built, numerically stable on \texttt{sm\_121}, and in the CUDA
  images; it lifted dense prefill $+42\%$ (Table~\ref{tab:flashattn}) and closed the gap from
  $0.55\times$, but this is the one axis Hanzo still trails. The residual $\approx 20\%$ is the
  \texttt{f32}$\leftrightarrow$\texttt{bf16} cast chain plus MMQ GEMM efficiency, the next two levers.
  \item \textbf{$35$B hybrid \emph{total} prefill on CUDA.} The GDN recurrence kernel itself now leads
  \texttt{llama.cpp} ($0.93$ vs $1.04$ ms/layer) and total prefill rose $354 \to 1166$ tok/s; the
  remaining gap on the full hybrid is a \texttt{bf16} tensor-core GEMM for the projections and experts,
  the next item.
  \item \textbf{Tiny-model decode dispatch overhead ($0.6$B HIP $0.69\times$).} Fixed per-token launch
  cost dominates a small weight stream; the decode graph (which already won on CUDA) and further
  op-fusion are the levers.
\end{enumerate}

Two walls we could not move, we proved were silicon rather than software: the RDNA3.5 \texttt{iu8}
WMMA rate (equal to \texttt{f16}, so quantization is a bandwidth play not a compute play on this APU)
and the decode matvec's $88$--$99\%$ bandwidth saturation. Naming a ceiling as physics is itself a
result; it tells the next engineer where \emph{not} to spend effort.

\section{Conclusion}
\label{sec:conclusion}

The bottom line is deliberately unembellished. Hanzo ML and Hanzo Engine are one composable,
multi-backend, multi-modal inference stack: one quantization core serving the full 1-bit-to-8-bit
GGUF zoo bit-exact across CUDA, ROCm, Metal, Vulkan, and CPU, with a constant-cost add-a-quant
contract (Proposition~\ref{prop:o1}); one engine spanning text, vision, audio, speech, image
generation, and embeddings, including a unified omni model. Measured against \texttt{llama.cpp} (and
thus Ollama) on the same silicon, Hanzo \emph{leads} on AMD prefill for every model, leads CUDA MoE
prefill ($1.87$--$2.1\times$) and the CUDA GDN recurrence kernel ($0.93$ vs $1.04$ ms/layer), leads on
decode where the regime allows (the $30$B MoE at $1.20\times$ HIP, the $0.6$B at $1.10\times$ CUDA, the
\texttt{IQ3\_XXS} i-quant at $1.13\times$ HIP), and reaches decode parity across all three backends. The
one axis it still trails is \emph{dense}-attention prefill on CUDA: a Blackwell flash-attention kernel,
now built and shipped, lifted it $+42\%$ and closed the gap to $0.77$--$0.84\times$ \texttt{llama.cpp}
but did not erase it, and the residual cast-and-GEMM levers are stated as engineering, not hidden. The
remaining frontiers, APU decode versus Vulkan and the $35$B hybrid total prefill, are likewise named
with the lever and forward motion measured. The campaign's
ethos is the paper's credibility: we profile the production path, gate every change bit-exact,
control the device map and the thermals, and report the disproven ideas as first-class negative
results. The data leads a lot; we let it carry, and we are precise where it does not.

\begin{thebibliography}{10}

\bibitem{hanzorocm}
Hanzo AI. \textit{Native ROCm Inference on a Consumer RDNA3.5 APU: Reaching llama.cpp Decode Parity
via a Unified 1-bit-to-Full Quantization Compute Core}. Hanzo AI white paper (companion), 2026.

\bibitem{llamacpp}
G. Gerganov et al. \textit{llama.cpp: LLM inference in C/C++}. 2023.

\bibitem{ggml}
G. Gerganov et al. \textit{ggml: Tensor library for machine learning}. 2023.

\bibitem{ollama}
Ollama. \textit{Ollama: Get up and running with large language models locally} (built on
\texttt{llama.cpp}/\texttt{ggml}). 2023.

\bibitem{qwen3}
Qwen Team. \textit{Qwen3 Technical Report}. Alibaba, 2025.

\bibitem{williams2009}
S. Williams, A. Waterman, and D. Patterson. Roofline: An Insightful Visual Performance Model for
Multicore Architectures. \textit{Communications of the ACM}, 52(4):65--76, 2009.

\bibitem{dao2022}
T. Dao et al. FlashAttention: Fast and Memory-Efficient Exact Attention with IO-Awareness.
\textit{NeurIPS}, 2022.

\bibitem{kwon2023}
W. Kwon et al. Efficient Memory Management for Large Language Model Serving with PagedAttention.
\textit{SOSP}, 2023.

\bibitem{gateddeltanet}
S. Yang et al. \textit{Gated Delta Networks: Improving Mamba2 with Delta Rule}. 2024.

\bibitem{rocm2024}
AMD. \textit{ROCm: Open Software Platform for GPU Compute}. Versions 7.1 and 7.13, 2026.

\end{thebibliography}

\end{document}
